% Original: PDF strana 23, gornji slajd.
\slajdid{02-s045}
\begin{frame}[label=02-s045]{Poređenje po broju ulaza kola}
\setlength{\parskip}{0pt}
% element: 02-s045:e01
\begin{columns}[T,onlytextwidth]
\begin{column}{.48\textwidth}
\Oznaka{Zbir proizvoda}
\[F=\bar D\bar B+D\bar C\bar A\]
Jedno dvoulazno I, jedno troulazno I i jedno dvoulazno ILI kolo.
\[2+3+2=\mathbf{7}\text{ ulaza}\]
\end{column}
\begin{column}{.48\textwidth}
\Oznaka{Proizvod zbirova}
\[F=(D+\bar B)(\bar D+\bar C)(\bar D+\bar A)\]
Tri dvoulazna ILI i jedno troulazno I kolo.
\[3\cdot2+3=\mathbf{9}\text{ ulaza}\]
\end{column}
\end{columns}
\par\vspace{3mm}
% element: 02-s045:e02
Za ovaj primer zbir proizvoda je povoljniji prema ukupnom broju ulaza I i ILI kola.
\par\vspace{2mm}
\Rezultat{\alert{\textbf{Unapred se ne zna koja je realizacija povoljnija.}} Treba izvesti \alert{\textbf{oba minimalna izraza}} — zbir proizvoda i proizvod zbirova — pa njihovim poređenjem zaključiti koja je realizacija minimalna prema izabranom kriterijumu.}
\end{frame}
